% 19 November 2026. Drafted for Edward Baker with OpenAI Codex assistance.
% Root settings unavailable; B03 Astra/ultra and E14 Sol6.1/ultra dispatches.
% Integration of the E14 proof; assessment status is recorded in Review07.
\section{Full-phase correction and the pointwise strip consumer}
\label{sec:uniform-correction}

The fixed-cell gain raises two different questions: whether the correction
can be controlled through the whole Riemann--Siegel phase range, and whether
the resulting approximation proves a positive final margin. We resolve
the first and compare complete pointwise allowances for the second.
The comparison retains an explicit unacquired arithmetic lower bound;
it proves no new zero-free region.

\subsection{One family and removable apparent poles}

Keep \eqref{cr:reference}, now with the natural index and the family
\begin{equation}
 \frac9{100}\le t\le\frac1{10},\qquad x\ge14000,\qquad |y|\le1,
 \quad W=\frac{x}{2}+\frac{\pi t}{8},\quad
 w=W+\frac{\ii y}{2},\quad
 N=\left\lfloor\sqrt{\frac W{2\pi}}\right\rfloor.
 \label{uc:family}
\end{equation}
For clarity distinguish the published value correction and its lift:
\begin{align}
 C_{\rm val,N}&=J_N A_N(W),& Q_{\rm val,N}&=P_N-C_{\rm val,N},\\
 C_{\rm hol,N}&=J_N A_N(w),& Q_{\rm hol,N}&=P_N-C_{\rm hol,N}.
 \label{uc:two-references}
\end{align}
Only the second reference with a \emph{fixed} index is holomorphic.
The natural-index patchwork is not asserted to be holomorphic across
cutoff seams. The value reference is sufficient for a pointwise
nonvanishing test, which takes no spatial derivative.

\begin{lemma}[Full-phase analytic bounds]
The functions $C_0$ and $C_0^\sharp$ in \eqref{cr:reference} extend to
entire even functions. On
$|\Re p|\le101/100$, $|\Im p|\le1/100$ they satisfy
\begin{equation}
 |C_0|<4,\qquad |C_0'|<20,\qquad |C_0''|<200,
 \label{uc:C-bounds}
\end{equation}
and the same bounds hold for the sharp function. Every natural phase
and its vertical continuation in \eqref{uc:family} lies in this patch.
\end{lemma}

\begin{proof}
At $p=j+1/2$, the exponential in the numerator is
$\ii(-1)^{j(j+1)/2}$, while
$\sqrt2\cos(\pi j/2+\pi/4)=(-1)^{j(j+1)/2}$.
The numerator therefore vanishes at every simple zero of the
denominator. Removal of these singularities proves entire extension;
evenness is inherited from the original quotient.

On the boundary of $|\Re p|\le5/4$, $|\Im p|\le1/4$, the denominator
has modulus at least $\sqrt2$: this follows on vertical sides from
$|\cos(a+\ii b)|^2=\cos^2a+\sinh^2b$, and on horizontal sides from
$2\sinh(\pi/4)>\sqrt2$. The numerator has modulus below
$e^{5\pi/16}+\sqrt2\cosh(\pi/8)<5$. Hence the quotient is below
$5/\sqrt2<4$, including the interior by the maximum principle.
Cauchy disks of radius $1/5$ centered on the smaller rectangle remain
inside the larger one, proving the derivative bounds.

On the vertical segment through $W\ge7000$,
$|p_N'|\le1/\sqrt{2\pi W}$. Its phase displacement from the real
phase is at most $1/(2\sqrt{2\pi W})<1/100$.
The real natural phase lies in $(-1,1]$, with $-1$ included by
one-sided limits. This proves the asserted full phase coverage, without
excluding neighborhoods of $\pm1/2$.
\end{proof}

\subsection{Derivative, mismatch and complete source payments}

\begin{proposition}[Uniform correction payments]
Throughout \eqref{uc:family}, with primes denoting derivatives in $w$,
\begin{equation}
 |A_N|<5W^{3/2},\qquad |A_N'|<9W,\qquad
 |A_N''|<40\sqrt W.
 \label{uc:A-bounds}
\end{equation}
For fixed $N$, the complete heat defect is
\begin{equation}
 R_{Q_{\rm hol,N}}=R_{P_N}-\frac{J_N A_N''}{4},\qquad
 |R_{Q_{\rm hol,N}}|\le |R_{P_N}|+10\sqrt W\,|J_N|.
 \label{uc:defect}
\end{equation}
The value-to-lift mismatch is bounded by
\begin{equation}
 |C_{\rm hol,N}-C_{\rm val,N}|\le M_N
 :=\frac92 W|y||J_N|,\qquad
 \frac{M_N}{L}\le\frac{9|y|}{\sqrt W}\le\frac9{\sqrt{7000}}<\frac19,
 \quad L=\frac{|J_N|}{2}\sqrt W\sqrt{W^2+1}.
 \label{uc:mismatch}
\end{equation}
\end{proposition}

\begin{proof}
Put $B=((w+\ii)C_0+(w-\ii)C_0^\sharp)/2$, so $A_N=\sqrt w B$.
Writing $d=(2\pi W)^{-1/2}$, the full product and chain rules give
\begin{align*}
 |B|&\le4(W+3/2),\\
 |B'|&\le4+20(W+3/2)d,\\
 |B''|&\le40d+(W+3/2)\{200d^2+10d/W\}.
\end{align*}
Use $|\sqrt w|\le\sqrt{W+1/2}$,
$|(\sqrt w)'|\le1/(2\sqrt W)$ and
$|(\sqrt w)''|\le1/(4W^{3/2})$.
Substitution into $A_N'=(\sqrt w)'B+\sqrt w B'$ and
$A_N''=(\sqrt w)''B+2(\sqrt w)'B'+\sqrt w B''$
gives \eqref{uc:A-bounds}; for explicit overpayments one can use
$\sqrt W>80$, $d<(5/12)/\sqrt W$,
$W+3/2<(1001/1000)W$ and
$\sqrt{W+1/2}<(1001/1000)\sqrt W$.
All terms, including $p_N''$, are retained in the displayed bounds.

The $A_N$ and $A_N'$ terms cancel in $\partial_t+\partial_{zz}$
exactly as in \eqref{cr:defect}, giving \eqref{uc:defect}.
Integration along the vertical segment gives the mismatch.
Finally $W^2+1\ge W^2$ and $7000>81^2$ prove the last comparison.
This is a triangle defect bound, with no assumed signed cancellation.
\end{proof}

Let $\mathcal V_N=S_N+L(\delta_-+\delta_+)$ denote the complete
retained source budget from \eqref{cr:source}, using the natural $N$ in
both summand-error sums and the full retained remainder. The source
estimate of \cite[Corollary 6.4]{Polymath2019} and
\eqref{uc:mismatch} give
\begin{align}
 |H-P_N|&\le V_{\rm om}:=\mathcal V_N+L,\notag\\
 |H-Q_{\rm val,N}|&\le V_{\rm val}:=\mathcal V_N,\notag\\
 |H-Q_{\rm hol,N}|&\le V_{\rm hol}:=\mathcal V_N+M_N.
 \label{uc:budgets}
\end{align}
In particular $V_{\rm om}-V_{\rm hol}>8L/9$. This saves a fraction
of the omitted leading payment, not of the entire error or final margin.

\subsection{A quantitative comparison of complete pointwise allowances}

\begin{theorem}[Retained-value improvement and the direct comparator]
On \eqref{uc:family} define
\[
 m_{\rm om}=|P_N|-V_{\rm om},\qquad
 m_{\rm val}=|Q_{\rm val,N}|-V_{\rm val},\qquad
 m_{\rm hol}=|Q_{\rm hol,N}|-V_{\rm hol}.
\]
These are lower bounds for $|H|$, and
\begin{equation}
 \boxed{m_{\rm val}>m_{\rm om}+\frac3{40}L,
 \qquad m_{\rm hol}\le m_{\rm val}.}
 \label{uc:complete-comparison}
\end{equation}
Both comparisons remain true after division by the same positive
carrier modulus. Neither asserts that $m_{\rm val}$ is positive.
\end{theorem}

\begin{proof}
For real $p\in[-1,1]$, \cite[(58)]{Polymath2019} gives $|C_0(p)|\le1/2$.
There is also the elementary bound
$0<\Re C_0(p)\le\cos(\pi/8)/2$. To prove it, put
$q=|p|$, $a=q-1/2$ and $b=(q+1/2)/2$. Cancellation in the real part gives
\[
 \Re C_0(p)=\frac{\sin(\pi ab)}{2\sin(\pi a)}.
\]
The expression is interpreted continuously at $a=0$.
For $u=\pi|a|\in[0,\pi/2]$ and $b\in[1/4,3/4]$,
$\sin(bu)/\sin u$ increases in both arguments. Its $u$ derivative has
the sign of $b\cot(bu)-\cot u$, which is positive because $v\cot v$
decreases on $(0,\pi/2]$. The maximum is therefore
$\sin(3\pi/8)=\cos(\pi/8)$.

Since $A_N(W)=\sqrt W\{W\Re C_0-\Im C_0\}$,
\[
 \frac{|C_{\rm val,N}|}{L}
 \le\frac{W\cos(\pi/8)+1}{\sqrt{W^2+1}}
 \le\cos(\pi/8)+\frac1W<\frac{37}{40}.
\]
For the last rational bound use $\cos(\pi/8)<231/250$ and
$1/W\le1/7000<1/1000$. The reverse triangle inequality now gives
$m_{\rm val}-m_{\rm om}\ge L-|C_{\rm val,N}|>3L/40$.
Finally $Q_{\rm hol,N}-Q_{\rm val,N}=C_{\rm val,N}-C_{\rm hol,N}$,
so $|Q_{\rm hol,N}|\le|Q_{\rm val,N}|+M_N$, proving the second claim.
\end{proof}

The first comparison retains correction information that the omitted
budget discards. The second gives that same information to the direct
consumer. It is restricted to exact pointwise reference values and a
separately paid vertical triangle allowance. It does not compare arbitrary
interval implementations, computational costs, correlated residuals,
counted-pair functionals or heat propagation. In particular it does not
withdraw the fixed-cell discriminant gain of Section~9.

\subsection{Cutoff seams and the remaining arithmetic input}

At $W=2\pi j^2$, equivalently $x=4\pi j^2-\pi t/4$, the old index is
$j-1$ and the new one is $j$. Put
$d_j=p_j(t,s_+)+p_j(t,s_-)$, as in \eqref{as:cutoff-jump}.
Since $J_{j-1}=-J_j$ and the real seam phases are $-1,+1$,
evenness gives the exact signed jumps
\begin{align}
 \Delta Q_{\rm val}&=d_j-2J_jA_j(W),\notag\\
 \Delta Q_{\rm hol}&=d_j-J_j\{A_j(w)+A_{j-1}(w)\}.
 \label{uc:seams}
\end{align}
The two profiles in the second line are not identified off the real axis.
Both one-sided source estimates apply at the seam by continuity; hence
$|\Delta Q|\le V_{\rm old}+V_{\rm new}$ for each type of reference,
with its own complete budgets. A change of reference transfers the error
by $E_{\rm new}=E_{\rm old}-\Delta Q$, including the jump payment.
Pointwise coverage uses the natural half-open cells and their endpoints.
Differentiating the floor, or treating the patchwork as a single heat
solution without seam terms, is not justified by these formulas.

The chosen final consumer is precisely \eqref{as:strip-final} with
\begin{equation}
 t_0=1/10,\quad y_0=1/5,\quad X=14000,\qquad
 x\ge14000+\sqrt{24/25},\quad 1/5\le y\le\sqrt{4/5}.
 \label{uc:consumer}
\end{equation}
The nonvanishing carrier $\mathcal B_{t_0}$ in
\eqref{as:normalizer} permits the normalized statement. What remains
unacquired is an independently proved lower function
\[
 \ell(x,y)\le
 \left|\frac{Q_{\rm val,N}(t_0,x+\ii y)}{\mathcal B_{t_0}(x+\ii y)}\right|,
 \qquad
 \ell(x,y)>\frac{\mathcal V_N(t_0,x+\ii y)}
 {|\mathcal B_{t_0}(x+\ii y)|}
\]
at every point of \eqref{uc:consumer}, with a complete finite-height/tail
split and all cutoff interfaces. A positive unscaled infimum over
unbounded height is not requested. Defining $\ell$ from the desired
positivity would not acquire it.

The family includes the finite interval
$[14000,14000+\sqrt{24/25})$, but acquires no positivity there.
It supplies neither the initial verification \eqref{as:strip-initial}
nor the barrier \eqref{as:strip-barrier} down to time zero. Only if
all three criterion inputs were established would the parameter
arithmetic yield $\Lambda\le1/10+1/50=3/25$; this is not a new bound.

The collision screens retain their precise failed premises. The
Euler-factor, weighted-square, finite-prime and reflected/positive-primitive
controls do not inherit the literal theta source budget. That failure
does not turn a small approximation error into a zero-sensitive margin:
at a zero, any valid complete allowance gives $|Q|-V\le0$.
The quadratic $H=Q=(z-X)^2-2(t-T)$ has zero approximation error and
still collides. Finite-source controls lack the required all-real anchor;
RH-conditional escaping controls keep their hypothesis and location
qualifiers; changed observation kernels pay their actual forcing.
Thus the complete comparison is useful, while final positivity remains
an arithmetic acquisition problem.
